\input{preamble/preamble.tex}

\geometry{margin=0.85in}
\AtBeginDocument{\small}

\newcommand{\ExamNameField}{\noindent\textbf{Name:}\ \rule{0.7\linewidth}{0.4pt}\par\medskip}

\begin{document}
	\section*{10th grade - Learning evidence 2.1: Analysis of Decisions (Practice)}\label{sec:grade10-exam2-1-practice}
	
	\subsection*{Assessment and evaluation}\label{sec:grade10-exam2-1-practice-evaluation}
	
	This exam evaluates the following criteria. Each criterion is evaluated across the three points below.
	\begin{itemize}
		\item C2 Interprets decision alternatives, events, consequences, and states of nature.
		\item C3 Builds a payoff table from a description of the problem.
		\item C4 Explains the maximax criterion for decision making without probabilities.
		\item C5 Summarizes the maximin criterion for decision making without probabilities.
		\item C1 Describes the way a problem can be formulated for optimal decision making.
	\end{itemize}
	A student passes a criterion if it is correctly activated in at least two of the three points.
	The overall exam result is binary (Passed/Not passed) and is determined by the set of criteria passed.
	
	\begin{enumerate}
		\item A campus bookstore must choose one of three inventory plans for a two-week launch: Plan A (expanded textbook stock),
		Plan B (balanced mix), or Plan C (lean stock). Demand for books can be high, medium, or low. Based on past launches, the probability of high demand is $0.40$, medium demand is $0.35$, and low demand is $0.25$. Profits are measured in hundreds of dollars. If demand is high, Plan A yields 88, Plan B yields 76, and Plan C yields 60. If demand is medium, Plan A yields 52, Plan B yields 58, and Plan C yields 50. If demand is low, Plan A yields $-12$, Plan B yields 28, and Plan C yields 36. The stated probabilities must be used for an expected value comparison as part of the decision analysis.
		
		\begin{itemize}
			\item Interpret the decision alternatives, events, consequences, and states of nature in this context.
			\item Build the payoff table from the description.
			\item Explain the maximax criterion without probabilities and state the recommended plan.
			\item Summarize the maximin criterion without probabilities and state the recommended plan.
			\item Describe and apply the expected value criterion using the stated probabilities, then interpret the result.
		\end{itemize}
		\item A regional courier company is choosing a fleet strategy for the next season: Strategy A (electric vans),
		Strategy B (hybrid vans), or Strategy C (diesel vans). Demand for deliveries can be strong, steady, or weak, and fuel costs can be low or high. The company estimates the demand probabilities as $0.25$ (strong), $0.55$ (steady), and $0.20$ (weak). Fuel costs are independent of demand, with probabilities $0.60$ (low) and $0.40$ (high). Revenue, in thousands of dollars, depends on the demand level and strategy. Under strong demand, revenue is 460 for Strategy A, 410 for Strategy B, and 360 for Strategy C. Under steady demand, revenue is 330 for Strategy A, 310 for Strategy B, and 280 for Strategy C. Under weak demand, revenue is 220 for Strategy A, 240 for Strategy B, and 230 for Strategy C. Operating costs, in thousands of dollars, depend on the fuel cost level and strategy. If fuel costs are low, operating costs are 170 for Strategy A, 200 for Strategy B, and 220 for Strategy C. If fuel costs are high, operating costs are 240 for Strategy A, 270 for Strategy B, and 300 for Strategy C. The company should treat each combined demand and fuel-cost outcome as a state of nature and use the stated probabilities to compute expected profits.
		
		\begin{itemize}
			\item Interpret the decision alternatives, events, consequences, and states of nature, including the combined states implied by independence.
			\item Build the payoff table by computing profits for each combined state using the revenue and cost descriptions.
			\item Explain the maximax criterion without probabilities and state the recommended strategy.
			\item Summarize the maximin criterion without probabilities and state the recommended strategy.
			\item Describe and apply the expected value criterion using the combined-state probabilities, then interpret the result.
		\end{itemize}
		\item A community recreation center must choose a membership model for the next year: Model A (standard monthly pass),
		Model B (premium pass with classes), or Model C (flex pass). Demand for memberships can be high, medium, or low, and staffing costs can be favorable or unfavorable. Demand probabilities are $0.30$ (high), $0.45$ (medium), and $0.25$ (low). Staffing costs are independent of demand, with probabilities $0.55$ (favorable) and $0.45$ (unfavorable). Expected membership counts depend on demand and model: under high demand the center expects 640 members with Model A, 500 with Model B, and 720 with Model C; under medium demand it expects 480 with Model A, 380 with Model B, and 540 with Model C; under low demand it expects 320 with Model A, 250 with Model B, and 380 with Model C. Membership fees are $36$ for Model A, $50$ for Model B, and $32$ for Model C per member per year. Variable staffing cost per member is $19$ when costs are favorable and $27$ when costs are unfavorable. Fixed annual costs are $6{,}800$ for Model A, $9{,}200$ for Model B, and $5{,}400$ for Model C. All monetary amounts in this point can be treated in dollars. Use the combined demand and staffing outcomes as states of nature and the stated probabilities to compute expected profits.
		
		\begin{itemize}
			\item Interpret the decision alternatives, events, consequences, and states of nature in this setting.
			\item Build the payoff table by calculating profit for each combined state using the narrative data.
			\item Explain the maximax criterion without probabilities and state the recommended model.
			\item Summarize the maximin criterion without probabilities and state the recommended model.
			\item Describe and apply the expected value criterion using the combined-state probabilities, then interpret the result.
		\end{itemize}
	\end{enumerate}
\end{document}
